% ===============================================================
\section{Introduction}\label{sec:intro}
% ===============================================================

\subsection{Background and Motivation}\label{subsec:background}

Retrieval-Augmented Generation (RAG) grounds language-model outputs in
external evidence, improving factual reliability without
retraining~\cite{lewis2020rag,guu2020realm}. Its most common form relies on
dense passage retrieval, embedding query and corpus into a shared vector
space and returning the most semantically similar chunks~\cite{karpukhin2020dpr}.
This is effective for factoid questions answered in a single passage, but it
treats the corpus as a flat collection of independent chunks and reduces
retrieval to nearest-neighbour search over isolated
spans~\cite{karpukhin2020dpr,trivedi2023interleaving}.

In practice, this assumption is overly simplistic for many real-world
question-answering settings. When the answer depends on relations among
entities, on document structure, or on chains of evidence spanning multiple
documents, dense retrieval cannot follow those connections because the
relations are not represented in the vector space~\cite{trivedi2023interleaving}.
Graph-based retrieval was introduced precisely to address this gap: by
organising the corpus as a knowledge graph, retrieval can traverse explicit,
typed relationships rather than relying solely on surface
similarity~\cite{edge2024graphrag,wang2024kgp,gutierrez2024hipporag}. Yet
graph traversal is more expensive than dense lookup, and applying it to every
query---including simple ones whose answer lies in a single article---incurs
unnecessary latency and computational cost. Vietnamese legal question
answering exposes this tension sharply: the same system must serve both
simple article-level definitions and multi-document questions involving
amendment, replacement, and cross-document
reference~\cite{chalkidis2020legalbert,zhong2020nlp_law,sovrano2020legalkg},
while Vietnamese-specific word segmentation and multi-syllable named entities
further complicate legal reference matching~\cite{nguyen2020phobert,vu2018vncorenlp}.

When confronted with such heterogeneous queries, a question naturally arises:
how to decide \emph{which} retrieval strategy a given query requires,
\emph{before} any evidence is retrieved. Existing Hybrid RAG systems sidestep
this question by applying dense and graph retrieval indiscriminately, paying
the cost of both regardless of query
complexity~\cite{edge2024graphrag,he2024lightrag,gao2023modularrag}. Recent
adaptive-routing methods show that selecting a strategy by query complexity
improves both quality and
efficiency~\cite{jeong2024adaptive,chen2023frugalgpt,ong2024routellm}, but
they route among \emph{model sizes} or \emph{retrieval depths} rather than
among distinct \emph{evidence mechanisms}; the decision of whether a query
needs relational traversal at all is left unaddressed. A separate and
safety-critical gap concerns ambiguity. When a user asks \textit{``Văn bản đó
còn hiệu lực không?''} (\textit{``Is that document still in effect?''})
without a usable referent in the conversation, forcing retrieval grounds the
answer in an arbitrary legal document and returns a confidently wrong
answer~\cite{min2020ambigqa,elgohary2019canard}. No prior system
simultaneously decides among dense retrieval, graph traversal, hybrid
reasoning, and clarification for Vietnamese legal question answering.

Motivated by the observation that the appropriate retrieval strategy is a
property of the query itself---and can be predicted cheaply from the query
and its conversational context---we propose a two-stage reasoning-aware
adaptive router that selects one of four evidence actions before retrieval. A
lightweight XGBoost classifier~\cite{chen2016xgboost} handles clear cases at
low cost; a selective LLM verifier is invoked only for uncertain, ambiguous,
or relation-heavy queries, following the cost-aware cascade
principle~\cite{chen2023frugalgpt,ong2024routellm}. A deterministic history
resolver precedes both stages to distinguish resolved, irrelevant, and
conflicting conversation references, so that clarification is triggered only
when the legal target is genuinely underspecified.

\subsection{Main Contributions}\label{subsec:contributions}

To address the heterogeneity of legal queries described above, we adopt a
design principle in which the retrieval strategy is selected by a two-stage
router operating upstream of retrieval, rather than by running all strategies
and merging their outputs. The first stage performs a fast, feature-based
classification of the query into a retrieval mechanism using an XGBoost model
trained on routing labels derived from evidence structure. Because this stage
is lightweight, it can serve as a real-time gateway that resolves the
majority of clear cases without invoking a large model. The second stage is
an LLM verifier triggered selectively, only when the first-stage decision is
uncertain, when the query carries strong relational signals, or when the
conversational referent cannot be resolved. This selective design keeps the
expected per-query cost low while reserving explicit reasoning for the cases
that need it. Crucially, the second stage does more than correct routing
decisions: by reasoning explicitly over query intent and legal constraints
before answering, it also functions as a chain-of-thought step that improves
the quality of the generated answer. Clarification is integrated as a
first-class routing action so that, when no retrieval strategy can safely
ground the answer, the system asks the user rather than guessing.

Our contributions are four-fold. \textbf{First}, we formulate
reasoning-aware retrieval routing for Vietnamese legal QA as a four-class
decision---dense retrieval, graph traversal, hybrid reasoning, and
clarification---made before evidence retrieval, unifying mechanism selection
and ambiguity handling in a single
framework~\cite{jeong2024adaptive,min2020ambigqa}. \textbf{Second}, we design
a two-stage router in which a lightweight XGBoost stage handles clear cases
and a selective LLM verifier is invoked only for uncertain, ambiguous, or
relation-heavy queries, and we show that this verifier acts as an explicit
chain-of-thought mechanism: the Two-stage Hybrid ($\text{F1}=0.640$)
outperforms an Oracle Router with perfect routing labels
($\text{F1}=0.563$), and a controlled \textit{Oracle+Stage-2} ablation that
holds routing at gold labels while enabling the verifier isolates a causal
$+0.074$~F1 contribution of second-stage reasoning ($0.637$ vs $0.563$),
demonstrating that the gain comes from reasoning during generation rather
than from route correction alone~\cite{chen2023frugalgpt}. \textbf{Third}, we introduce a deterministic
history resolver that classifies a conversational referent as resolved,
irrelevant, or conflicting, enabling the router to unblock retrieval when a
referent is available and to force clarification when it is
absent~\cite{elgohary2019canard}. \textbf{Fourth}, through experiments on a
strict $600$-query benchmark, a $234$-query ambiguity benchmark, and a
$160$-query conversation stress test, we validate (a)~end-to-end answer
quality against fixed-strategy and oracle baselines, (b)~routing accuracy
against rule-based and learned baselines, (c)~clarification behaviour, and
(d)~the cost--quality trade-off of selective verification.

To the best of our knowledge, this is the first system to treat
clarification as a first-class routing action alongside dense, graph, and
hybrid retrieval in a Hybrid GraphRAG pipeline for Vietnamese legal question
answering~\cite{sovrano2020legalkg,zhong2020nlp_law}.

\subsection{Organisation}\label{subsec:organisation}

Section~\ref{sec:related_work} reviews related work.
Section~\ref{sec:methodology} defines the routing problem, dataset, and
pipeline. Section~\ref{sec:routing_framework} details the routing framework.
Sections~\ref{sec:eval} and~\ref{sec:results} present the experimental setup
and results. Section~\ref{sec:limitations} discusses limitations before
Section~\ref{sec:conclusion} concludes.
